%% kapa: exam=1 round=35 year=2024 period=1 form=A subject=civil_law subject_ko=민법 seq=24 number=24
%% source: Q-Net 감정평가사 기출문제 (cst00309) artlSeq=5213229 file=raw/1008/5213229/1교시 문제지 원본.pdf
%% conversion: kordoc 4.24.51 -> tools/gen1.py (auto); review=proofread-sonnet; flags=none
\begin{question}{24}
甲이 乙소유 X도자기에 관해 무단으로 丙에게 질권을 설정해 주었고, 丙은 질권의 선의취득을 주장하고 있다. 이에 관한 설명으로 옳지 않은 것은? (다툼이 있으면 판례에 따름)
\begin{choices}
  \item 丙은 평온ㆍ공연하게 X의 점유를 취득하였어야 한다.
  \item 丙은 甲이 소유자가 아니라는 사실에 대하여 그 자신이 선의이고 무과실이라는 사실을 증명하여야 한다.
  \item 丙이 甲과 질권설정계약을 체결할 당시 선의였다면 질물의 인도를 받을 때 악의라도 丙의 선의취득은 인정된다.
  \item 丙이 X에 대하여 甲이 직접점유를 취득하는 형태로 점유를 취득한 경우, 丙의 선의취득은 인정되지 아니한다.
  \item 만약 甲이 미성년자임을 이유로 丙과의 질권설정계약을 취소하면 丙은 선의취득을 할 수 없다.
\end{choices}
\end{question}
\answer{3}
